\documentclass[10pt,a4paper]{amsart}
\usepackage[margin=19mm]{geometry}
\usepackage{amsmath,amssymb,mathtools}
\usepackage[scr=rsfs,cal=euler]{mathalfa}
\usepackage{xfrac}
\usepackage[unicode,hidelinks,bookmarks=false]{hyperref}
\newcommand{\ie}{i.e.\ }
\newcommand{\cf}{cf.\ }
\newcommand{\resp}{respectively\ }
\newcommand{\Subsection}{\subsection}
\providecommand{\nobreakdash}{\nobreak\hbox{-}}
\setlength{\emergencystretch}{2em}
\multlinegap0pt
\usepackage{amssymb}
\newtheorem{lemma}{Lemma}
\newtheorem{theorem}{Theorem}
\title{On Lacunarity and Sidon-Type Properties for Subsystems of Characters of Compact Abelian Groups}
\author{Anna Kazakova}
\date{Source: arXiv:2605.01572v2 (CC BY 4.0)}
\begin{document}
\maketitle

\noindent\emph{Source and licence.} On Lacunarity and Sidon-Type Properties for Subsystems of Characters of Compact Abelian Groups. Version arXiv:2605.01572v2 (CC BY 4.0), distributed by arXiv under the Creative Commons Attribution 4.0 International licence, http://creativecommons.org/licenses/by/4.0/, as recorded in the arXiv licence metadata for this exact version and shown on the abstract page https://arxiv.org/abs/2605.01572v2. Full licence notice: \texttt{licenses/arxiv-2605.01572-CC-BY-4.0.txt}.\par\smallskip

\noindent\textit{Editorial scope.} Counted unit: the article's theorem th-0 (source0.tex lines 348--365). The article's complete original proof of it is delivered with it (source0.tex lines 366--438, one \texttt{proof} environment of 73 lines). No mathematics is written, repaired or re-derived: the statement and the proof are the article's own lines, in the article's own notation, and no retained line is substituted or re-flowed.  Retained uncounted as the source-local context the proof needs: lines 93--95; lines 312--323; lines 316--322; lines 335--337; lines 339--344.  Nothing was omitted from any retained span, so the package carries no omission record.  No retained line contains a citation, so the wrapper carries no bibliography and no bibliography entry.  No retained line contains a figure environment, an image-inclusion command or drawing code, so no image is delivered and the package declares no asset dependency.  Editorial formatting only, in addition: an AMS \texttt{amsart} preamble that restates the article's own declarations and macros used by the retained lines, namely the environments \texttt{lemma}, \texttt{theorem} on separate counters, together with the standard packages those lines need.  The retained environments print in this document as Lemma 1, Theorem 1 in order of appearance, whereas the article prints the same items with the numbering of its own full document.  The complete native source of the article as distributed in the arXiv e-print for this exact version is bundled unchanged as \texttt{sources/199645/source0.tex}.
	\begin{equation} \label{m_k}
		\text{where} \; m_k = \inf\{m\in \mathbb{N} \cup \{+\infty\} \colon \gamma_k^m =1\}.
	\end{equation}	

	\begin{lemma} \label{muy}
Let $R_i$ denote the generalized Rademacher functions with base $2d+1$. Then, for each fixed $y \in [0,1]$, there exists a measure $\rho_y$ such that for any fixed indices $k_1 < \cdots < k_s$ and $\alpha_j \in \{1, \ldots, d_{k_j}\}$, the corresponding coefficients of $\rho_y$ are given by the following formula:
		\begin{equation}  \label{muycoef}
			\widehat{\rho_y}( \gamma_{k_1}^{\alpha_1} \cdot \ldots \cdot \gamma_{k_s}^{\alpha_s}) = \frac{ R_{k_1}^{\alpha_1^{\prime}}(y) \cdot \ldots \cdot R_{k_s}^{\alpha_s^{\prime}}(y)}{(2d)^{s}},   \end{equation}
		\begin{equation} \label{alpha'}
			\text{where} \quad \alpha'_i = \alpha'_i(\alpha_i, k_i) := 
			\begin{cases} 
				2d+1-\alpha, & \text{if } \exists\, \alpha < \alpha_i : \gamma_{k_i}^{\alpha} = \overline{\gamma_{k_i}^{\alpha_i}}, \\ 
				\alpha_i, & \text{otherwise}.
			\end{cases} 
		\end{equation}
	\end{lemma}	

		\begin{equation} 
			\text{where} \quad \alpha'_i = \alpha'_i(\alpha_i, k_i) := 
			\begin{cases} 
				2d+1-\alpha, & \text{if } \exists\, \alpha < \alpha_i : \gamma_{k_i}^{\alpha} = \overline{\gamma_{k_i}^{\alpha_i}}, \\ 
				\alpha_i, & \text{otherwise}.
			\end{cases} 
		\end{equation}

	\begin{equation} \label{rho_x}
		\rho_x =  weak^*\lim_{n \rightarrow\infty} \prod_{i=0}^{n} \left( 1 + \frac{1}{2d}\left( \sideset{}{'}{\sum}_{k} \gamma_i^{-k}(x) R_i^{k}  + \sideset{}{'}{\sum}_{k} \overline{\gamma_i^{-k}(x) R_i^{k}}\right) \right),
	\end{equation}

	\begin{equation} \label{rhox} 
		\begin{split}
			& \quad \widehat{\rho_x}(R_{k_1}^{\alpha_1^{\prime}}(y) \cdot \ldots \cdot R_{k_s}^{\alpha_s^{\prime}}(y)) = \frac{\gamma_{k_1}^{-\alpha_1} \cdot \ldots \cdot \gamma_{k_s}^{-\alpha_s}}{(2d)^{s}}, \\
			& \quad \widehat{\rho_x}(R_{k_1}^{-\alpha_1^{\prime}}(y) \cdot \ldots \cdot R_{k_s}^{-\alpha_s^{\prime}}(y)) = \overline{\widehat{\rho_x}(R_{k_1}^{\alpha_1^{\prime}}(y) \cdot \ldots \cdot R_{k_s}^{\alpha_s^{\prime}}(y))} = \frac{\gamma_{k_1}^{\alpha_1} \cdot \ldots \cdot \gamma_{k_s}^{\alpha_s}}{(2d)^{s}}, 
		\end{split} 
	\end{equation}

	\begin{theorem} \label{th-0}
Let $\{\gamma_{k}\}_{k=0}^{\infty}$ be a $d$-dissociated system of characters of a compact abelian group, and let $m_k$ be given by \eqref{m_k}. Then, for every $q \ge 1$, including $q = \infty$, and for every polynomial of the form
		\[
		Q^{(s)} := \sum_{0\le k_1 < \ldots < k_s \le N} \; \sum_{\substack{\alpha_1+\ldots+\alpha_s = d \\ 1 \le \alpha_i < m_{k_i}}} C_{k_1, \ldots, k_s}^{\alpha_1, \ldots, \alpha_s} \gamma_{k_1}^{\alpha_1}(x) \cdot \ldots \cdot \gamma_{k_s}^{\alpha_s}(x), \qquad s \le d
		\]
		there exist polynomials 
		\begin{equation} \label{Q-s_R}
			\begin{split}
				& Q^{(s)}_{R, -}(y) := \sum_{0\le k_1 < \ldots < k_s \le N} \; \sum_{\substack{\alpha_1+\ldots+\alpha_s = d \\ 1 \le \alpha_i < m_{k_i}}} C_{k_1, \ldots, k_s}^{\alpha_1, \ldots, \alpha_s} R_{k_1}^{-\alpha^{\prime}_1}(y) \cdot \ldots \cdot R_{k_s}^{-\alpha^{\prime}_s}(y), \\
				& Q^{(s)}_{R, +}(y) := \sum_{0\le k_1 < \ldots < k_s \le N} \; \sum_{\substack{\alpha_1+\ldots+\alpha_s = d \\ 1 \le \alpha_i < m_{k_i}}} C_{k_1, \ldots, k_s}^{\alpha_1, \ldots, \alpha_s} R_{k_1}^{\alpha^{\prime}_1}(y) \cdot \ldots \cdot R_{k_s}^{\alpha^{\prime}_s}(y),
			\end{split}
		\end{equation}
		such that:
		\begin{equation} \label{th-0-main}
			\frac{1}{(2d)^{2d}} \|Q^{(s)}_{R, +}\|_q	\le \|Q^{(s)}\|_q \le (2d)^{2d} \|Q^{(s)}_{R, -}\|_q,
		\end{equation}
		where $\{R_{k}\}_{k=0}^{\infty}$ are Rademacher functions with base $2d+1$ and $\alpha_j'$ are defined by the formula \eqref{alpha'}.
	\end{theorem}

	\begin{proof}
		Consider the following polynomials from $x\in G$ and $y\in [0,1)$:
		\[
		\begin{split}
			& Q^{(s)}_{-}(x,y) = \sum_{0\le k_1 < \ldots < k_s \le N} \; \sum_{\substack{\alpha_1+\ldots+\alpha_s = d \\ 1 \le \alpha_i < m_{k_i}}} C_{k_1, \ldots, k_s}^{\alpha_1, \ldots, \alpha_s} R_{k_1}^{-\alpha^{\prime}_1}(y) \cdot \ldots \cdot R_{k_s}^{-\alpha^{\prime}_s}(y) \gamma_{k_1}^{\alpha_1}(x) \cdot \ldots \cdot \gamma_{k_s}^{\alpha_s}(x),\\
			& Q^{(s)}_{+}(x,y) = \sum_{0\le k_1 < \ldots < k_s \le N} \; \sum_{\substack{\alpha_1+\ldots+\alpha_s = d \\ 1 \le \alpha_i < m_{k_i}}} C_{k_1, \ldots, k_s}^{\alpha_1, \ldots, \alpha_s} R_{k_1}^{\alpha^{\prime}_1}(y) \cdot \ldots \cdot R_{k_s}^{\alpha^{\prime}_s}(y) \gamma_{k_1}^{\alpha_1}(x) \cdot \ldots \cdot \gamma_{k_s}^{\alpha_s}(x),
		\end{split}
		\]
where the tuple $(\alpha_1', \ldots, \alpha_s')$ is constructed from $(\alpha_1, \ldots, \alpha_s)$ and $(k_1, \ldots, k_s)$ via formula \eqref{alpha'}. For a fixed $y \in [0,1)$, consider the measure $\rho_y$ defined in Lemma \ref{muy}. Applying \eqref{muycoef}, we obtain: 
		\[
		Q^{(s)}_{-}(x,y) * \rho_y(x) = \int_{G} Q^{(s)}_{-} (x\dot-z, y) d \rho_y(z) 
		\]
		\[
		=\sum_{0\le k_1 < \ldots < k_s \le N} \; \sum_{\substack{\alpha_1+\ldots+\alpha_s = d \\ 1 \le \alpha_i < m_{k_i}}} C_{k_1, \ldots, k_s}^{\alpha_1, \ldots, \alpha_s} R_{k_1}^{-\alpha^{\prime}_1}(y) \cdot \ldots \cdot R_{k_s}^{-\alpha^{\prime}_s}(y) \cdot  \gamma_{k_1}^{\alpha_1}(x) \cdot \ldots \cdot \gamma_{k_s}^{\alpha_s}(x) \cdot \widehat{\rho_y}(\gamma_{k_1}^{\alpha_1} \cdot \ldots \cdot \gamma_{k_s}^{\alpha_s})
		\]
		\[
		= \frac{1}{(2d)^s} \sum_{0\le k_1 < \ldots < k_s \le N} \; \sum_{\substack{\alpha_1+\ldots+\alpha_s = d \\ 1 \le \alpha_i < m_{k_i}}} C_{k_1, \ldots, k_s}^{\alpha_1, \ldots, \alpha_s} \gamma_{k_1}^{\alpha_1}(x) \cdot \ldots \cdot \gamma_{k_s}^{\alpha_s}(x) = \frac{Q^{(s)}(x)}{(2d)^s}.
		\]
		Then, by Young's inequality, 
		\begin{equation} \label{Qy}
			\|Q^{(s)}(x)\|_{L_q(x)} = (2d)^{s}\|Q^{(s)}_{-}(x,y)* \rho_y\|_{L_q(x)} \le (2d)^{d}\|Q^{(s)}_{-}(x,y)\|_{L_q(x)}. 
		\end{equation}
		Taking the $L_q$-norm with respect to $y$ of both sides of \eqref{Qy}, we obtain:
		\begin{equation} \label{Qy2}
			\|Q^{(s)} (x)\|_{L_q(x)} \le (2d)^d \|Q^{(s)}_{-}(x,y)\|_{L_q(x)\times L_q(y)}. 
		\end{equation}
		
	Consider the polynomial $Q^{(s)}_{R, -}$, defined by the formula \eqref{Q-s_R}, and consider for a fixed $x\in G$ the measure $\rho_x$, defined by the formula \eqref{rho_x}. Calculate $Q^{(s)}_{R, -}(y)*\rho_x$, taking into account \eqref{rhox},
		\[
		Q^{(s)}_{R, -}(y) * \rho_x(y) = \int_0^1 Q^{(s)}_{R, -} (y\dot-z) d \rho_x(z) 
		\]
		\[
		=\sum_{0\le k_1 < \ldots < k_s \le N} \; \sum_{\substack{\alpha_1+\ldots+\alpha_s = d \\ 1 \le \alpha_i < m_{k_i}}} C_{k_1, \ldots, k_s}^{\alpha_1, \ldots, \alpha_s} R_{k_1}^{-\alpha^{\prime}_1}(y) \cdot \ldots \cdot R_{k_s}^{-\alpha^{\prime}_s}(y) \cdot \widehat{\rho_x}(R_{k_1}^{-\alpha^{\prime}_1}(y) \cdot \ldots \cdot R_{k_s}^{-\alpha^{\prime}_s}(y))
		\]
		\[
		= \frac{1}{(2d)^s} \sum_{0\le k_1 < \ldots < k_s \le N} \; \sum_{\substack{\alpha_1+\ldots+\alpha_s = d \\ 1 \le \alpha_i < m_{k_i}}} C_{k_1, \ldots, k_s}^{\alpha_1, \ldots, \alpha_s} R_{k_1}^{-\alpha^{\prime}_1}(y) \cdot \ldots \cdot R_{k_s}^{-\alpha^{\prime}_s}(y) \gamma_{k_1}^{\alpha_1}(x) \cdot \ldots \cdot \gamma_{k_s}^{\alpha_s}(x) = \frac{Q^{(s)}_{-}(x,y)}{(2d)^s}.
		\]
		Then, by Young's inequality, 
		\begin{equation} \label{Qx}
			\|Q^{(s)}_{-}(x,y)\|_{L_q(y)} = (2d)^{s}\|Q^{(s)}_{R,-}(y)* \rho_x\|_{L_q(y)} \le (2d)^{d}\|Q^{(s)}_{R,-}(y)\|_{L_q(y)}. 
		\end{equation}
		Taking the $L_q$-norm with respect to $x$ of both sides of \eqref{Qx}, we obtain:
		\begin{equation} \label{Qx2}
			\|Q^{(s)}_{-}(x,y)\|_{L_q(x)\times L_q(y)} \le (2d)^{d}\|Q^{(s)}_{R,-}(y)\|_{L_q(y)}.
		\end{equation}
		From \eqref{Qx2} and \eqref{Qy2} we get:
		\[
		\|Q^{(s)} (x)\|_{L_q(x)} \le (2d)^{2d} \|Q^{(s)}_{R,-}(y)\|_{L_q(y)}.
		\]
		Thus, the right inequality in \eqref{th-0-main} is proved.
		
        The left-hand inequality is proved similarly. For a fixed $y \in [0,1]$, using \eqref{muycoef}, we obtain:  
		\[
		Q^{(s)} * \rho_y = \int_{G} Q^{(s)} (x\dot-z) d\rho_y(z) = \frac{Q^{(s)}_{+}(x,y)}{(2d)^s}.
		\]
Applying Young's inequality and taking the $L_q$-norm with respect to $y$ (as in \eqref{Qy} and \eqref{Qy2}), we obtain:
		\begin{equation} \label{Qy2-l}
			\frac{1}{(2d)^d}\|Q^{(s)}_{+}(x,y)\|_{L_q(x)\times L_q(y)} \le \|Q^{(s)} (x)\|_{L_q(x)}.
		\end{equation}
		Taking into account \eqref{rhox} and \eqref{Q-s_R}, we obtain:
		\[
		Q^{(s)}_{+}(x,y) * \rho_x = \int_0^1 Q^{(s)}_{+} (x, y \dot-z)d\rho_x(z) = \frac{Q^{(s)}_{R, +}(y)}{(2d)^s}.
		\]
		Applying Young's inequality and taking the $L_q$-norm with respect to $x$ (as in \eqref{Qx} and \eqref{Qx2}), we obtain:
		\begin{equation} \label{Qx2-l}
			\frac{1}{(2d)^{d}} \|Q^{(s)}_{R,+}(y)\|_{L_q(y)}  \le  \|Q^{(s)}_{+}(x,y)\|_{L_q(x)\times L_q(y)}.
		\end{equation}
		From \eqref{Qy2-l} and \eqref{Qx2-l} we get the left inequality in \eqref{th-0-main}:
		\[
		\frac{1}{(2d)^{2d}} \|Q^{(s)}_{R,+}(y)\|_{L_q(y)} \le \|Q^{(s)} (x)\|_{L_q(x)}.
		\]	 
		This completes the proof.
	\end{proof}

\end{document}
